\section{Residual Stress and the Dipole}
\label{sec:residual}

\begin{definition}[Residual stress]
\derived\ Residual stress is the mechanical expression of the Angular Shortfall:
\[
\sigma_{\mathrm{res}}=\kappa\cdot\Delta\theta,\qquad\kappa\neq 0.
\]
\end{definition}

\begin{proposition}[Residual-stress dipole]
\derived\ Because residual stress must be dynamically active, $\kappa$ cannot
vanish. The only two continuous possibilities are opposite signs:
\[
P>0\quad\text{(tension)},\qquad p<0\quad\text{(pressure)}.
\]
Neither may be zero; both act on the same geometric medium. The unreachable null
separates them but does not forbid imperfect exchange.
\end{proposition}

The Kinetic optical mode (Section~\ref{sec:optics}) is the ordinary living regime
in which both signs remain active---the default state of residual organisation.
